\section{推理时提示与令牌预算}

\subsection{扩大 token 预算以暴露推理链}
Chen 以 Arc AGI puzzle 为例展示推理时 token 预算的曲线：传统模型在 budget 约 20 美元时止步 25\% 以下，而 OpenAI O3 在设定 20~1000 美元预算、配合示例提示与执行 trace 的前提下，可把准确率推到 87.5\%。此处幻灯片（图 \ref{fig:token-budget}）清楚展示推理精度随 token 预算拉长腰部的非线性趋势。

\begin{figure}[H]
\centering
\includegraphics[width=0.92\textwidth]{frame-token-budget.jpg}
\caption{Arc AGI 的推理时曲线与 OpenAI O3 的 token-cost 布局。\protect\footnotemark}
\label{fig:token-budget}
\end{figure}
\footnotetext{视频画面时间区间：00:08:56--00:10:20，使用 O3 曲线说明 token 扩展可达到 human-level accuracy。}

\begin{importantbox}{推理时曲线的三重启示}
1) 当任务复杂度升高时，token 预算需要从几十推到上千才能看到显著提升；2) 增加预算必须配合示例提示与 stepwise scoring，否则仅靠更长的 prompt 无法抬升准确率；3) 高预算模式适用于竞赛级别或验证任务，不适合原生低延迟对话。
\end{importantbox}

\subsection{提示设计的约束}
提示需要两条能力：一是鼓励模型生成长 Chain-of-Thought，让多轮合理推理序列自然浮现；二是将任务难度 signal 内置在 prompt 中，使模型自动调节生成深度。Chen 在 44 分 31 秒明确指出，“提示只能整出更长推理的前提是，被提示的模型知道自己面对的是难题”。

\begin{knowledgebox}{提示设计双标准}
1. 长 CoT：使用分步模板、问答对或多例示范，让模型持续追踪当前状态；2. 难度标记：把 problem difficulty、context length、expected depth 显示在 prompt 中，让模型在 inference-time 自适应调整 depth。
\end{knowledgebox}

\subsection{成本可观测性与预算控制}
扩展 token 预算要参考真实成本：从 1 美元、20 美元到 1,000 美元每题的区别，不只是钱，还带来 latency 与算力消耗。陈列的现实风险在于“高预算让任务不可解释，也可能吞噬微服务资源”。

\begin{warningbox}{预算跑飞的风险}
一旦推理时预算拉到上千美元，就必须限制在高价值子任务，否则会让调用频率踏空、系统吞吐率下降甚至引发平衡失控。建议提前设定 sampling 缓存策略，把最昂贵的 search 留给 verified-critical path。
\end{warningbox}

\subsection{本章小结}
推理时提示需要在“鼓励长 CoT”与“任务难度感知”之间找到平衡，同时把 token 预算、成本与 latency 指标纳入监控。Arc AGI 曲线图提醒我们，token 数量扩展有效但成本逐步饱和，必须让提示、验证与预算协调运作。

